The peak of an empirical learning curve can depend on how a solver defines numerical rank. We compare minimum-norm least squares, cross-validated ridge and principal components regression on artist-disjoint music data and an exact-feature-grouped materials holdout. Each comparison uses 100 training permutations at 15 sample sizes. Ridge has the lowest mean squared error of the three estimators at every size. On materials, one training subset accounts for 99.92\% of minimum-norm loss at the mean curve's peak, $n=73$ with 81 features. In a post-hoc recomputation, imposing a relative singular-value cutoff of $10^{-12}$ changes the fitted rank in just five of 1,500 subsets. The mean peak moves to $n=83$, while the median peak remains at $n=81$. Exact rank certificates on the recorded decimal features show that the worst subset has centered rank 71. The default fit retained all 71 directions. Thus a meaningful change in the peak can result from discarding an algebraically nonzero direction. The findings quantify how numerical regularization and the choice of loss summary affect the reading of a finite-benchmark learning curve.
